\documentclass[11pt,a4paper]{article}
\usepackage[utf8]{inputenc}
\usepackage[T1]{fontenc}
\usepackage{amsmath,amsfonts,amssymb}
\usepackage{graphicx}
\usepackage{geometry}
\usepackage{xcolor}
\usepackage{hyperref}
\usepackage{fancyhdr}
\usepackage{titlesec}
\usepackage{enumitem}
\usepackage{booktabs}
\usepackage{array}
\usepackage{mdframed}
\usepackage{listings}
\geometry{margin=2.5cm, top=3cm, bottom=3cm, headheight=15pt}

% Colors
\definecolor{mainblue}{RGB}{0,74,127}
\definecolor{accentblue}{RGB}{0,120,200}
\definecolor{lightgray}{RGB}{245,245,245}
\definecolor{codegray}{RGB}{240,240,240}
\definecolor{darkgray}{RGB}{80,80,80}

% Hyperref config
\hypersetup{
    colorlinks=true,
    linkcolor=mainblue,
    urlcolor=accentblue,
    citecolor=mainblue,
    pdftitle={Rapport des TPs - Réalité Augmentée From Scratch},
    pdfauthor={DIFFO KENNE GARNEL},
}

% Section title formatting
\titleformat{\section}{\large\bfseries\color{mainblue}}{}{0em}{\thesection\quad}[\vspace{-0.3em}\rule{\textwidth}{0.8pt}]
\titleformat{\subsection}{\normalsize\bfseries\color{accentblue}}{}{0em}{\thesubsection\quad}
\titleformat{\subsubsection}{\normalsize\bfseries\color{darkgray}}{}{0em}{\thesubsubsection\quad}

\titlespacing*{\section}{0pt}{1.5em}{0.8em}
\titlespacing*{\subsection}{0pt}{1em}{0.4em}
\titlespacing*{\subsubsection}{0pt}{0.8em}{0.3em}

% Header/footer
\pagestyle{fancy}
\fancyhf{}
\fancyhead[L]{\small\color{darkgray}Réalité Augmentée From Scratch}
\fancyhead[R]{\small\color{darkgray}DIFFO KENNE GARNEL — 24P816}
\fancyfoot[C]{\small\color{darkgray}\thepage}
\renewcommand{\headrulewidth}{0.4pt}
\renewcommand{\footrulewidth}{0pt}
\renewcommand{\headrule}{\color{mainblue}\hrule width\headwidth height\headrulewidth}

% Listings style
\lstset{
    backgroundcolor=\color{codegray},
    basicstyle=\ttfamily\small,
    frame=single,
    framerule=0.4pt,
    rulecolor=\color{accentblue},
    breaklines=true,
    numbers=none,
    xleftmargin=8pt,
    xrightmargin=4pt,
    aboveskip=6pt,
    belowskip=6pt,
}

% Framed environment for key formulas
\mdfdefinestyle{formulabox}{
    backgroundcolor=lightgray,
    linecolor=accentblue,
    linewidth=1pt,
    leftmargin=0pt,
    rightmargin=0pt,
    innerleftmargin=10pt,
    innerrightmargin=10pt,
    innertopmargin=8pt,
    innerbottommargin=8pt,
}

\begin{document}
\thispagestyle{empty}

% ============================================================
%  PAGE DE GARDE
% ============================================================
\begin{titlepage}
    \centering
    \vspace*{1.5cm}

    % Top rule
    {\color{mainblue}\rule{\textwidth}{2pt}}\\[0.4cm]

    {\fontsize{26}{32}\selectfont\bfseries\color{mainblue} RAPPORT DES TPs}\\[0.3cm]
    {\large\bfseries\color{accentblue} RÉALITÉ AUGMENTÉE FROM SCRATCH}\\[0.5cm]

    {\color{mainblue}\rule{\textwidth}{2pt}}\\[1.2cm]

    {\large\itshape\color{darkgray} Stabilité numérique, algèbre linéaire et pipeline graphique}\\[2cm]

    % Info block
    \begin{mdframed}[style=formulabox]
    \renewcommand{\arraystretch}{1.5}
    \begin{tabular}{@{}ll}
        \textbf{UE} & Coder pour la VR, l'XR et l'AR 2 \\
        \textbf{Étudiant} & DIFFO KENNE GARNEL \\
        \textbf{Matricule} & 24P816 \\
        \textbf{Établissement} & ENSPY \\
        \textbf{Date} & 03 Avril 2026 \\
    \end{tabular}
    \end{mdframed}

    \vfill

    {\small\color{darkgray} Module NkMath — Semaines 1 à 4}\\[0.5cm]
    {\color{mainblue}\rule{\textwidth}{1pt}}
\end{titlepage}

\newpage
\tableofcontents
\newpage

% ============================================================
\section{Introduction}
% ============================================================

Ce rapport présente l'ensemble des travaux pratiques réalisés dans le cadre du cours \textit{Réalité Augmentée from Scratch}, couvrant les quatre premières semaines du module \textbf{NkMath}. L'objectif principal est de construire une bibliothèque mathématique complète en C++17 from scratch, sans dépendances tierces, en se concentrant sur les thèmes suivants :

\begin{itemize}[leftmargin=2em, itemsep=4pt]
    \item \textbf{Semaine 1} — Stabilité numérique et représentation IEEE 754 des flottants
    \item \textbf{Semaine 2} — Structures vectorielles (Vec2d, Vec3d, Vec4d) et algèbre linéaire de base
    \item \textbf{Semaine 3} — Matrices 4$\times$4, transformations affines et inversion numérique
    \item \textbf{Semaine 4} — Quaternions, interpolation sphérique (SLERP) et animation
\end{itemize}

Chaque section détaille les concepts théoriques, les implémentations pratiques et les tests unitaires correspondants. L'ensemble est couvert par plus de \textbf{40 tests} validant la robustesse numérique et les cas limites.

\newpage

% ============================================================
\section{Précision Numérique}
% ============================================================

\subsection{Représentation IEEE 754 des flottants}

Un nombre flottant IEEE 754 simple précision (32 bits) est structuré comme suit :

\begin{itemize}[itemsep=2pt]
    \item \textbf{Bit 31} : Signe (0 = positif, 1 = négatif)
    \item \textbf{Bits 30–23} : Exposant biaisé (8 bits, biais $= 127$)
    \item \textbf{Bits 22–0} : Mantisse (23 bits, avec bit implicite 1)
\end{itemize}

La valeur représentée est :
\begin{equation}
(-1)^{s} \times 2^{(e-127)} \times \left(1 + m \times 2^{-23}\right)
\end{equation}

où $s$ est le bit de signe, $e$ l'exposant biaisé, et $m$ la mantisse.

\subsubsection{Pourquoi 0.1f n'est pas représentable exactement}

$0.1_{10} = 0.000\overline{1100}_2$ en binaire. La mantisse de 23 bits ne peut représenter que les 23 premiers bits après la virgule, conduisant à l'approximation :
\[
0.1 \approx 0.100000001490116119384765625
\]
Cette erreur d'arrondi est la source de nombreux problèmes numériques dans les calculs AR.

\subsubsection{Valeurs spéciales IEEE 754}

\begin{center}
\renewcommand{\arraystretch}{1.4}
\begin{tabular}{llll}
\toprule
\textbf{Valeur} & \textbf{Exposant} & \textbf{Mantisse} & \textbf{Remarque} \\
\midrule
$\pm 0$      & 0   & 0        & Zéros signés \\
$\pm \infty$ & 255 & 0        & Résultat d'overflow \\
NaN          & 255 & $\neq 0$ & Contagieux, tester avec \texttt{std::isnan()} \\
Subnormal    & 0   & $\neq 0$ & Très petits nombres \\
\bottomrule
\end{tabular}
\end{center}

\subsection{Epsilon machine et ULP}

L'epsilon machine $\varepsilon$ est la plus petite valeur telle que $1.0 + \varepsilon \neq 1.0$ :
\[
\varepsilon_{\text{float}} \approx 1.19 \times 10^{-7}, \qquad \varepsilon_{\text{double}} \approx 2.22 \times 10^{-16}
\]

L'ULP (\textit{Unit in the Last Place}) mesure l'espacement entre flottants consécutifs. Il varie avec la magnitude :

\begin{itemize}[itemsep=2pt]
    \item Près de $1.0$ : ULP $\approx 1.19 \times 10^{-7}$
    \item Près de $1000.0$ : ULP $\approx 1.22 \times 10^{-4}$
\end{itemize}

\subsection{Erreurs d'arrondi et algorithmes stables}

\subsubsection{Cancellation catastrophique}

La soustraction de deux nombres presque égaux perd des bits significatifs :
\[
a = 1\,234\,567.89,\quad b = 1\,234\,567.00 \implies a - b = 0.89 \approx 0.875
\]

\subsubsection{Sommation de Kahan}

L'algorithme compense les erreurs accumulées à chaque étape :
\begin{align}
y &= x_i - c \\
t &= s + y \\
c &= (t - s) - y \\
s &= t
\end{align}
où $c$ est le terme de compensation.

\subsubsection{Variance : méthode naïve vs Welford}

La méthode naïve :
\[
\sigma^2 = \frac{\displaystyle\sum x_i^2 - \left(\sum x_i\right)^2/n}{n-1}
\]
est instable numériquement (peut produire une variance négative).

La méthode de Welford maintient la précision en une seule passe :
\begin{align}
\delta   &= x_i - \bar{x}_{i-1} \\
\bar{x}_i &= \bar{x}_{i-1} + \delta/i \\
M2_i     &= M2_{i-1} + \delta\,(x_i - \bar{x}_i)
\end{align}

\subsection{Résultats expérimentaux}

\begin{center}
\renewcommand{\arraystretch}{1.3}
\begin{tabular}{lcc}
\toprule
\textbf{Méthode} & \textbf{Résultat} ($10^6$ additions de 0.1f) & \textbf{Erreur} \\
\midrule
\texttt{std::accumulate} & 100\,011.55 & $+11.55$ \\
Kahan & 100\,000.00 & $\approx 0$ \\
\bottomrule
\end{tabular}
\end{center}

Sur le jeu $\{10^8, 10^8, 1.0, 2.0\}$, la méthode naïve peut retourner une variance négative, tandis que Welford reste stable.

\newpage

% ============================================================
\section{Vecteurs et Algèbre Linéaire}
% ============================================================

\subsection{Vec2d : Vecteurs 2D}

Implémentation avec opérateurs arithmétiques complets, accès par index et fonctions utilitaires.

\subsubsection{Produit scalaire}
\[
\mathbf{a} \cdot \mathbf{b} = a_x b_x + a_y b_y
\]

\subsubsection{Produit vectoriel 2D (scalaire)}
\[
\mathbf{a} \times \mathbf{b} = a_x b_y - a_y b_x
\]
Donne l'aire orientée du parallélogramme défini par $\mathbf{a}$ et $\mathbf{b}$.

\subsubsection{Normalisation}
\[
\hat{\mathbf{u}} = \frac{\mathbf{u}}{\|\mathbf{u}\|}, \qquad \|\mathbf{u}\| = \sqrt{u_x^2 + u_y^2}
\]

\subsection{Vec3d : Vecteurs 3D}

\subsubsection{Produit vectoriel 3D}
\[
\mathbf{a} \times \mathbf{b} =
\begin{pmatrix}
a_y b_z - a_z b_y \\
a_z b_x - a_x b_z \\
a_x b_y - a_y b_x
\end{pmatrix}
\]
Suit la règle de la main droite.

\subsubsection{Orthogonalisation de Gram-Schmidt}

Construction d'une base orthonormée $\{\mathbf{u}_1, \mathbf{u}_2, \mathbf{u}_3\}$ à partir de trois vecteurs $\mathbf{a}, \mathbf{b}, \mathbf{c}$ :
\begin{align}
\mathbf{u}_1 &= \frac{\mathbf{a}}{\|\mathbf{a}\|} \\
\mathbf{v}_1 &= \mathbf{b} - \operatorname{proj}_{\mathbf{u}_1}(\mathbf{b}), \quad \mathbf{u}_2 = \frac{\mathbf{v}_1}{\|\mathbf{v}_1\|} \\
\mathbf{v}_2 &= \mathbf{c} - \operatorname{proj}_{\mathbf{u}_1}(\mathbf{c}) - \operatorname{proj}_{\mathbf{u}_2}(\mathbf{c}), \quad \mathbf{u}_3 = \frac{\mathbf{v}_2}{\|\mathbf{v}_2\|}
\end{align}
avec $\operatorname{proj}_{\mathbf{u}}(\mathbf{v}) = (\mathbf{v} \cdot \mathbf{u})\,\mathbf{u}$.

\subsubsection{Rejet orthogonal}
\[
\operatorname{rej}_{\mathbf{u}}(\mathbf{v}) = \mathbf{v} - \operatorname{proj}_{\mathbf{u}}(\mathbf{v})
\]

\subsection{Vec4d : Coordonnées homogènes}

Utilisées pour les transformations affines ; le quatrième composant $w$ permet d'encoder les translations dans une multiplication matricielle.

\subsubsection{Projection perspective simple}

Pour un point 3D $(X, Y, Z)$, la projection 2D est :
\[
u = f\,\frac{X}{Z}, \qquad v = f\,\frac{Y}{Z}
\]
avec $f$ le facteur de champ de vision.

\newpage

% ============================================================
\section{Matrices et Transformations}
% ============================================================

\subsection{Matrices 4$\times$4 homogènes}

Les matrices $4 \times 4$ représentent les transformations affines 3D sous la forme :
\[
M = \begin{pmatrix} \mathbf{R} & \mathbf{t} \\ \mathbf{0}^T & 1 \end{pmatrix}
\]
où $\mathbf{R} \in \mathbb{R}^{3\times3}$ est la rotation et $\mathbf{t} \in \mathbb{R}^3$ la translation.

\subsubsection{Composition TRS}

Les transformations Translation–Rotation–Scale se composent dans l'ordre :
\[
M = T \cdot R \cdot S
\]

\subsubsection{Décomposition TRS}

Extraction des paramètres depuis $M$ :
\begin{itemize}[itemsep=2pt]
    \item \textbf{Translation} : colonne 4 de $M$
    \item \textbf{Scale} : normes des colonnes de la sous-matrice $3\times3$
    \item \textbf{Rotation} : colonnes normalisées de la sous-matrice $3\times3$
\end{itemize}

\subsubsection{Inversion numérique de matrice}

Utilisation de l'algorithme de Gauss-Jordan avec \textbf{pivot partiel} pour assurer la stabilité. Testé sur matrices aléatoires et cas singuliers.

\subsubsection{Caméra LookAt}

Construction de la matrice de vue depuis un triplet $(\text{eye}, \text{target}, \text{up})$ :
\begin{align}
\mathbf{z} &= \operatorname{normalize}(\text{eye} - \text{target}) \\
\mathbf{x} &= \operatorname{normalize}(\text{up} \times \mathbf{z}) \\
\mathbf{y} &= \mathbf{z} \times \mathbf{x}
\end{align}

\subsubsection{Matrice de projection perspective}

\[
P = \begin{pmatrix}
\dfrac{f}{\text{aspect}} & 0 & 0 & 0 \\[8pt]
0 & f & 0 & 0 \\[6pt]
0 & 0 & \dfrac{\text{far}+\text{near}}{\text{near}-\text{far}} & \dfrac{2\cdot\text{far}\cdot\text{near}}{\text{near}-\text{far}} \\[8pt]
0 & 0 & -1 & 0
\end{pmatrix}
\]
avec $f = \dfrac{1}{\tan(\theta/2)}$ et $\theta$ le FOV vertical.

\newpage

% ============================================================
\section{Rasterisation Logicielle}
% ============================================================

\subsection{Pipeline graphique}

Le pipeline implémenté suit quatre étapes :

\begin{enumerate}[itemsep=4pt]
    \item \textbf{Transformation modèle} : espace objet $\to$ espace monde
    \item \textbf{Transformation vue} : espace monde $\to$ espace caméra
    \item \textbf{Projection} : espace caméra $\to$ NDC
    \item \textbf{Rasterisation} : NDC $\to$ pixels écran
\end{enumerate}

La transformation complète d'un point local vers l'écran s'écrit :
\[
\mathbf{p}_{\text{écran}} = P \cdot V \cdot M \cdot \mathbf{p}_{\text{local}}
\]

\subsection{Composants implémentés}

\subsubsection{Cube unitaire}

Définition des 8 sommets et 12 arêtes d'un cube centré en $(0, 0, 0)$.

\subsubsection{Algorithme de tracé de lignes}

Utilisation de l'algorithme de \textbf{Bresenham} pour rasteriser les arêtes entre sommets projetés.

\subsubsection{Format de sortie PPM}

Sauvegarde au format binaire P6 :
\begin{lstlisting}
P6
width height
255
<data RGB>
\end{lstlisting}

\subsubsection{Animation par rotation}

Rotation autour de l'axe Y :
\[
R_y(\theta) = \begin{pmatrix}
\cos\theta & 0 & \sin\theta & 0 \\
0          & 1 & 0          & 0 \\
-\sin\theta& 0 & \cos\theta & 0 \\
0          & 0 & 0          & 1
\end{pmatrix}
\]
Génération de \textbf{10 frames} pour simuler une rotation complète à $360°$.

\newpage

% ============================================================
\section{Quaternions et Animation}
% ============================================================

\subsection{Représentation des quaternions}

Un quaternion représente une rotation 3D :
\[
\mathbf{q} = w + x\,\mathbf{i} + y\,\mathbf{j} + z\,\mathbf{k}
\]
avec la contrainte d'unitarité $w^2 + x^2 + y^2 + z^2 = 1$.

\subsubsection{Rotation d'un vecteur}
\[
\mathbf{v}' = \mathbf{q}\,\mathbf{v}\,\mathbf{q}^{-1}
\]

\subsubsection{Conversion quaternion $\to$ matrice $3\times3$}

\[
R = \begin{pmatrix}
1-2y^2-2z^2 & 2xy-2wz     & 2xz+2wy \\
2xy+2wz     & 1-2x^2-2z^2 & 2yz-2wx \\
2xz-2wy     & 2yz+2wx     & 1-2x^2-2y^2
\end{pmatrix}
\]

\subsubsection{Inverse et identité}

Pour un quaternion unitaire : $\mathbf{q}^{-1} = \bar{\mathbf{q}} = (w,\,-x,\,-y,\,-z)$.

Vérification : $\mathbf{q} \cdot \mathbf{q}^{-1} = 1$.

La stabilité numérique du cycle $\mathbf{q} \to R \to \mathbf{q}$ a été vérifiée sur \textbf{50 cas aléatoires}.

\subsection{Interpolation sphérique (SLERP)}

\subsubsection{LERP vs SLERP}

\begin{itemize}[itemsep=3pt]
    \item \textbf{LERP} : interpolation linéaire dans l'espace des quaternions — vitesse angulaire \textit{non constante}
    \item \textbf{SLERP} : interpolation sur le grand cercle de $S^3$ — vitesse angulaire \textit{constante}
\end{itemize}

\subsubsection{Formule SLERP}

\[
\operatorname{slerp}(\mathbf{q}_1, \mathbf{q}_2, t) = \frac{\sin\!\bigl((1-t)\theta\bigr)}{\sin\theta}\,\mathbf{q}_1 + \frac{\sin(t\theta)}{\sin\theta}\,\mathbf{q}_2
\]
avec $\cos\theta = \mathbf{q}_1 \cdot \mathbf{q}_2$.

\subsubsection{Animation}

Interpolation entre deux orientations sur \textbf{60 frames}, avec rendu du cube en rotation et comparaison visuelle LERP/SLERP.

\newpage

% ============================================================
\section{Tests et Validation}
% ============================================================

Le fichier \texttt{TestsDesTPs.cpp} contient 11 suites de tests couvrant l'ensemble des fonctionnalités :

\begin{center}
\renewcommand{\arraystretch}{1.4}
\begin{tabular}{cll}
\toprule
\textbf{Suite} & \textbf{Thème} & \textbf{Description} \\
\midrule
TP1  & Flottants IEEE 754        & Inspection bit-à-bit (\texttt{inspectFloat}) \\
TP2  & Précision                 & Somme Kahan, variance Welford, epsilon machine \\
TP3  & Utilitaires flottants     & 33 tests : isFiniteValid, nearlyZero, approxEq \\
TP4  & Vec2d                     & 20 tests : dot, cross 2D, normalisation, accès \\
TP5  & Vec3d / Gram-Schmidt      & 16 tests : cross 3D, ortho. sur 10 triplets aléatoires \\
TP6  & Projection Vec4d          & Rendu d'un cube avec sommets projetés \\
TP7  & Mat4d / Inversion         & 10 tests : matrices aléatoires et singulières \\
TP8  & Rasterisation             & 10 frames d'animation du cube en rotation \\
TP9  & Décomposition TRS         & 20 tests sur matrices aléatoires \\
TP10 & Quaternions               & 50 tests $\mathbf{q}\to R\to\mathbf{q}$, 50 tests d'inverse \\
TP11 & Animation SLERP           & 60 frames SLERP et LERP pour comparaison \\
\bottomrule
\end{tabular}
\end{center}

\vspace{0.5em}
Au total, \textbf{plus de 40 assertions} couvrent les cas nominaux, les cas limites et la robustesse numérique.

\newpage

% ============================================================
\section{Conclusion}
% ============================================================

Ce projet a permis de construire les fondations mathématiques essentielles d'un moteur de réalité augmentée from scratch. Les principales réalisations sont :

\begin{itemize}[itemsep=5pt]
    \item \textbf{Précision numérique} — Analyse approfondie de l'IEEE 754 et implémentation des algorithmes stables Kahan et Welford
    \item \textbf{Algèbre vectorielle} — Structures Vec2d, Vec3d, Vec4d avec algèbre linéaire complète (dot, cross, Gram-Schmidt)
    \item \textbf{Matrices 4$\times$4} — Inversion numérique par Gauss-Jordan, composition et décomposition TRS
    \item \textbf{Pipeline graphique} — Rasterisation logicielle avec projection perspective, algorithme de Bresenham et export PPM
    \item \textbf{Quaternions} — Représentation compacte des rotations avec interpolation SLERP pour une animation fluide
\end{itemize}

L'ensemble est validé par plus de 40 tests unitaires couvrant les cas limites et la robustesse numérique. La bibliothèque \textbf{NkMath} constitue une base solide pour les modules suivants : SVD, traitement d'image et calibration caméra.

\end{document}
